\documentclass[12pt]{article}

\usepackage{fontawesome}
\usepackage{hyperref}
\usepackage{xurl}
\usepackage[tagged, highstructure]{accessibility}
\usepackage{bookmark}
\usepackage{enumitem}

\hypersetup{
    colorlinks=false,
    pdfborder={0 0 0},
    pdftitle={Introduction to Data - Datasets and TED Talks},
    pdfauthor={Adrianna Holden-Gouveia},
    pdfsubject={Advanced Data - Lab Assignment},
    pdflang={en-US},
    pdfkeywords={datasets, data quality, TED talks, data analysis}
}

% Configure PDF bookmarks for navigation
\bookmarksetup{
    numbered,
    open,
}

% Configure list spacing for better accessibility
\setlist{nosep}



\title{Introduction to Data}
\author{
        Adrianna Holden-Gouveia \\
        Website: \href{https://aholdengouveia.name}{aholdengouveia.name}\\ 
        \date{\vspace{-5ex}}
        %Email: \href{mailto:admin@aholdengouveia.name}{admin@aholdengouveia.name} \\
%        \faLinkedin{: aholdengouveia} \\
        \faGithub {: aholdengouveia} \\
        %\faTwitter {: aholdengouveia} \\
        }

%S\date{\today}


\begin{document}    

\maketitle

%\begin{abstract}
%This is a template for Linux Administration Lab work
%\end{abstract}
%\tableofcontents


\section*{Accessibility Notice}
This document is also available in HTML format at:

Link: \href{https://aholdengouveia.name/AdvData/labs/week1.html}{\textbf{\nolinkurl{https://aholdengouveia.name/AdvData/labs/week1.html}}}

The HTML version provides enhanced accessibility features including keyboard navigation, screen reader support, responsive design, dark mode support, and high contrast options.

\section*{Request}
\noindent\textbf{Request:} DATA-2614\\
\textbf{Requested by:} Head of the data team\\
\textbf{Data:} The data set you choose in this request

\section*{Background}
You've just started on the data team at a new company. Your first job is to find a data set the team can use for the rest of the term, check that it's reliable and that a company is allowed to use it, and recommend one. You decide what kind of company you work for, based on the data set you choose. As part of your first week of training, you'll also review a talk about how data is used.

\section*{Learning Outcomes}
By the end of this lab, you will be able to:
\begin{itemize}
    \item Find data sets from reliable sources
    \item Profile a data set, including its source, license, size, and quality
    \item Choose a data set for a business purpose, and justify your choice
    \item Summarize and review a talk about how data is used
\end{itemize}

\section*{Lab Instructions}
The data set you choose in this lab is the one you will use for the rest of the term, so choose carefully. You can't choose a data set about books, since books are my example data set throughout the term.

\subsection*{Data Profiles}
\begin{enumerate}
    \item Find 3 data sets. You may use your own resources or the ones on my Resources page (linked under Resources). Check that each source is reliable using the CRAAP test (linked under Resources).
    \item Make a data profile, called a data card, for each data set. Follow the format of the sample data card (linked under Resources). Each data card must include:
    \begin{itemize}
        \item The name, source with a link, and creator or owner, and a sentence on why the source is reliable
        \item The license, and whether a company could use the data commercially
        \item When it was published, and how often it is updated
        \item The format and size (number of rows and columns)
        \item The key columns and their data types
        \item How many values are missing in the columns you care about most
        \item Quality notes: anything messy, out of date, or hard to use
        \item At least 2 pros and 2 cons
        \item A type of company that would use this data, and a decision it would help them make
    \end{itemize}
\end{enumerate}

\subsection*{Choose Your Data Set}
Pick the one data set you will use for the rest of the term. In later labs, you'll load a sample of 100 to 1,000 rows of it into a database using my CSV to Database Converter (linked under Resources), so it's fine if the full data set is much bigger. Your data set needs to be available as a CSV file, or as a spreadsheet you can save as a CSV file. In a paragraph, explain why you chose it over the other two, using what you found in your data cards.

\subsection*{TED Talk Review}
Watch one of the following TED talks (or read the TED article) and write a 1 page review that answers these 5 questions:
\begin{enumerate}
    \item What is the main point about data?
    \item What evidence or examples are used to support it?
    \item Who benefits from the way data is used, and who could be harmed?
    \item Do you agree with the speaker? Why or why not?
    \item How does this connect to the data set you chose, or to working with data in a job?
\end{enumerate}
\begin{itemize}
    \item How data brokers sell your identity, Madhumita Murgia Link: \href{https://www.ted.com/talks/madhumita_murgia_how_data_brokers_sell_your_identity/transcript?language=en}{\textbf{\nolinkurl{https://www.ted.com/talks/madhumita_murgia_how_data_brokers_sell_your_identity/transcript?language=en}}}
    \item How to use data to make a hit TV show, Sebastian Wernicke Link: \href{https://www.ted.com/talks/sebastian_wernicke_how_to_use_data_to_make_a_hit_tv_show/transcript?language=en}{\textbf{\nolinkurl{https://www.ted.com/talks/sebastian_wernicke_how_to_use_data_to_make_a_hit_tv_show/transcript?language=en}}}
    \item How tech companies deceive you into giving up your data and privacy, Finn Lützow-Holm Myrstad Link: \href{https://www.ted.com/talks/finn_lutzow_holm_myrstad_how_tech_companies_deceive_you_into_giving_up_your_data_and_privacy/transcript?language=en}{\textbf{\nolinkurl{https://www.ted.com/talks/finn_lutzow_holm_myrstad_how_tech_companies_deceive_you_into_giving_up_your_data_and_privacy/transcript?language=en}}}
    \item Let my dataset change your mindset, Hans Rosling Link: \href{https://www.ted.com/talks/hans_rosling_let_my_dataset_change_your_mindset/transcript?language=en}{\textbf{\nolinkurl{https://www.ted.com/talks/hans_rosling_let_my_dataset_change_your_mindset/transcript?language=en}}}
    \item How bad data keeps us from good AI, Mainak Mazumdar Link: \href{https://www.ted.com/talks/mainak_mazumdar_how_bad_data_keeps_us_from_good_ai/transcript}{\textbf{\nolinkurl{https://www.ted.com/talks/mainak_mazumdar_how_bad_data_keeps_us_from_good_ai/transcript}}}
    \item How open government data creates smarter societies (TED article) Link: \href{https://ideas.ted.com/how-open-government-data-creates-smarter-societies/}{\textbf{\nolinkurl{https://ideas.ted.com/how-open-government-data-creates-smarter-societies/}}}
\end{itemize}

\section*{Acceptance Criteria}
This request is done when:
\begin{itemize}
    \item 3 data cards, each with all 9 items from the checklist
    \item Each data card links to its data set
    \item The license is checked on every data card, including whether a company can use the data commercially
    \item Your chosen data set isn't about books
    \item Your paragraph explains your choice using what you found in your data cards
    \item Your TED talk review is about 1 page and answers all 5 questions
\end{itemize}

\section*{Deliverables}
One file that includes:
\begin{itemize}
    \item Your 3 data cards, each with a link to its data set
    \item Your paragraph explaining which data set you chose and why
    \item Your TED talk review
\end{itemize}

Audience: your manager, the head of the data team.

\section*{Resources}
References, a video, a PowerPoint and some notes are available at my website
Link: \href{https://www.aholdengouveia.name/AdvData/IntroData.html}{\textbf{\nolinkurl{https://www.aholdengouveia.name/AdvData/IntroData.html}}}

Places to find data sets, on my Resources page
Link: \href{https://www.aholdengouveia.name/AdvData/Resources.html}{\textbf{\nolinkurl{https://www.aholdengouveia.name/AdvData/Resources.html}}}

Sample data card, showing the format to follow
Link: \href{https://www.aholdengouveia.name/AdvData/labs/SampleDataCard.txt}{\textbf{\nolinkurl{https://www.aholdengouveia.name/AdvData/labs/SampleDataCard.txt}}}

The CRAAP test, for checking your sources
Link: \href{https://web.archive.org/web/20250814000638/https://guides.lib.uchicago.edu/c.php?g=1241077&p=9082343}{\textbf{\nolinkurl{https://web.archive.org/web/20250814000638/https://guides.lib.uchicago.edu/c.php?g=1241077&p=9082343}}}

My guide to evaluating sources
Link: \href{https://www.aholdengouveia.name/Howtos/Sourceeval/sourceEvaluation.html}{\textbf{\nolinkurl{https://www.aholdengouveia.name/Howtos/Sourceeval/sourceEvaluation.html}}}

My CSV to Database Converter, which makes a database file from your CSV file
Link: \href{https://www.aholdengouveia.name/AdvData/labs/csv-to-database.html}{\textbf{\nolinkurl{https://www.aholdengouveia.name/AdvData/labs/csv-to-database.html}}}

\end{document}